\section{Dados e curadoria}\label{sec:sv_dados}

A base Adult~\cite{uci:adult} foi extraída do censo dos EUA de 1994 por Barry Becker e doada à UCI por Becker e Kohavi~\cite{kohavi:96}. A versão do Kaggle~\cite{kaggle:adult} tem \svLinhasBrutas{} registros, um por pessoa, com 14 atributos e o rótulo \texttt{income}, que diz se a renda anual passa de 50 mil dólares (\svPositivosBrutos\% passam). O arquivo é conferido pelo seu \emph{hash} SHA-256, para garantir que todos os experimentos partem do mesmo arquivo. A Tabela~\ref{tab:sv_atributos} lista os 10 atributos usados, com os nomes deste trabalho.

\begin{table}[ht]
\centering
\caption{Atributos usados: nome no arquivo original, nome neste trabalho e valores, após a curadoria.}
\label{tab:sv_atributos}
\footnotesize
\begin{tabular}{lll}
\toprule
Coluna original & Nome neste trabalho & Valores \\
\midrule
\texttt{age} & idade & 17 a 90 anos \\
\texttt{education.num} & anos de estudo & 1 a 16 \\
\texttt{capital.gain} & ganho de capital & 0 a \svGanhoMax{} dólares \\
\texttt{capital.loss} & perda de capital & 0 a 4.356 dólares \\
\texttt{hours.per.week} & horas por semana & 1 a 99 \\
\texttt{workclass} & classe de trabalho & 7 categorias \\
\texttt{marital.status} & estado civil & 7 categorias \\
\texttt{occupation} & ocupação & 14 categorias \\
\texttt{relationship} & relacionamento & 6 categorias \\
\texttt{native.country} & país de origem & 41 países \\
\texttt{income} & renda acima de 50 mil (rótulo) & 1 = \texttt{>50K}, 0 = \texttt{<=50K} \\
\bottomrule
\end{tabular}
\end{table}

\textbf{Linhas.} \svDuplicadas{} linhas eram cópias exatas de outras, inclusive no peso amostral, e foram removidas. Outras \svAusentes{} tinham valor ausente, marcado com ``?'', quase sempre na classe de trabalho e na ocupação ao mesmo tempo, ou no país. Nessas linhas, \svPositivosAusentes\% ganham mais de 50 mil, contra \svPositivos\% no resto; como são 7\% da base e preencher a ocupação inventaria um dado num dos atributos mais informativos, foram removidas. Restam \svLinhas{} pessoas, \svPositivos\% com renda acima de 50 mil.

\textbf{Atributos.} O \texttt{fnlwgt} é um peso amostral, quantas pessoas da população cada registro representa, e não uma característica da pessoa; sua correlação com o rótulo é de \svCorrFnlwgt, e ele foi removido. O \texttt{education} é a escolaridade em texto e corresponde, nível a nível, ao \texttt{education.num}; ficou só o número. Raça e sexo foram mantidos no arquivo, mas ficaram fora dos atributos de todos os modelos. Foi uma decisão do grupo: o objetivo do trabalho é estudar a SVM, e não queríamos um modelo que usasse raça ou sexo para prever a renda de alguém. Isso não garante, porém, que o modelo ignore essas características (Seção~\ref{sec:limitacoes}).

\textbf{Perfis repetidos e erro inevitável.} Sem o peso amostral, pessoas diferentes passam a ter os mesmos 10 atributos: \svRepetemPerfil{} linhas repetem o perfil de outra. \svPerfisConflito{} desses perfis aparecem com as duas rendas, em \svLinhasConflito{} linhas (\svPctConflito\% da base). Qualquer modelo dá a mesma resposta a pessoas com os mesmos dados; o melhor possível é responder, em cada perfil, a renda mais comum, o que deixa no mínimo \svErrosIrred{} erros (\svIrredPct\% da base). Nenhum modelo pode, portanto, passar de \svTeto\% de acurácia.
